\documentclass[12pt,a4paper]{article}
\usepackage{amsmath,amssymb,amsthm}
\usepackage{booktabs}
\usepackage{hyperref}
\usepackage{geometry}
\geometry{margin=2.5cm}
\usepackage{enumitem}

\newtheorem{theorem}{Theorem}[section]
\newtheorem{lemma}[theorem]{Lemma}
\newtheorem{corollary}[theorem]{Corollary}
\newtheorem{proposition}[theorem]{Proposition}
\theoremstyle{definition}
\newtheorem{definition}[theorem]{Definition}
\theoremstyle{remark}
\newtheorem{remark}[theorem]{Remark}

\title{The Born Rule from Recognition Science:\\
$P = |\psi|^2$ from $J$-Cost Phase Invariance}

\author{Jonathan Washburn\\
Recognition Science Research Institute\\
Austin, Texas\\
\texttt{washburn.jonathan@gmail.com}}

\date{March 2026}

\begin{document}
\maketitle

\begin{abstract}
We derive the Born rule $P(i) = |\psi_i|^2$ from Recognition
Science (RS) by two independent routes.  First, the RS path-weight
derivation: the path action $C[\gamma] = \int J(r(t))\,dt$ yields
weights $w = e^{-C}$; the amplitude--cost bridge $C = 2A$ gives
$|\psi|^2 = e^{-C}$, recovering the Born rule directly from the
cost functional.  Second, a uniqueness theorem: the probability
function must satisfy five structural requirements---%
(i)~non-negativity, (ii)~normalization, (iii)~$J$-cost phase
independence, (iv)~multiplicativity from path-action
additivity, and (v)~continuity.  We prove that the unique
function satisfying all five is $P(c_i) = |c_i|^2$.
Alternative rules $P = |c|^k$ with $k \neq 2$ are excluded.
The interference identity follows automatically.  The Born rule
is not postulated but is a theorem of the cost-geometric
framework.  All structural results are derived from the RS axioms~\cite{RS_Axioms}.
\end{abstract}

\section{Introduction}
\label{sec:intro}

The Born rule $P(i) = |\psi_i|^2$ is the link between the quantum
formalism and experimental outcomes~\cite{Born1926}.  In standard
quantum mechanics it is postulated.  Why the probability is
$|\psi|^2$ rather than $|\psi|$, $|\psi|^4$, or some other
function of the amplitude is one of the deepest questions in
quantum foundations~\cite{vN1932}.

Gleason's theorem (1957) derives $P = |\psi|^2$ from the
properties of frame functions on Hilbert spaces of dimension
$\geq 3$~\cite{Gleason1957}.  In Many-Worlds, Zurek's
``envariance'' argument derives the Born rule from entanglement
symmetry~\cite{Zurek2005}.  RS~\cite{RS_Axioms} provides two independent routes
to $P = |\psi|^2$:
\begin{enumerate}[label=(\alph*),nosep]
  \item A \emph{path-weight derivation}: the path action
  $C[\gamma] = \int J(r(t))\,dt$ yields weights
  $w = e^{-C}$; the amplitude--cost bridge $C = 2A$ gives
  $|\psi|^2 = e^{-C}$, recovering the Born rule.
  \item A \emph{uniqueness derivation}: non-negativity,
  normalization, phase independence, multiplicativity
  over independent events, and continuity jointly force
  $P = |c|^2$ as the unique probability function.
\end{enumerate}
We present both in this paper.  Route (a) is the RS-native
derivation; route (b) is a structural uniqueness theorem that
serves as an independent cross-check.

\section{Cost Functional Foundation}
\label{sec:cost}

RS derives all physical law from a single cost functional~\cite{RS_Axioms}.
The three axioms are:
\begin{enumerate}[label=(A\arabic*),nosep]
  \item $J(1) = 0$ \hfill\textit{(Normalization)}
  \item $J(xy) + J(x/y) = 2J(x)J(y) + 2J(x) + 2J(y)$ for all $x, y > 0$
  \hfill\textit{(Recognition Composition Law)}
  \item $J''_{\log}(0) = 1$, where $J_{\log}(t) := J(e^t)$
  \hfill\textit{(Calibration)}
\end{enumerate}

\begin{theorem}[Cost Uniqueness~\cite{RS_Axioms}]
\label{thm:J}
The unique continuous solution to (A1)--(A3) is
$J(x) = \tfrac{1}{2}(x + x^{-1}) - 1$, satisfying $J(x) \geq 0$
with equality iff $x = 1$.
\end{theorem}

\begin{proof}[Proof sketch]
Setting $H(t) = 1 + J(e^t)$ reduces (A2) to the d'Alembert equation
$H(s+t) + H(s-t) = 2H(s)H(t)$.  With $H(0)=1$ and $H''(0)=1$,
the Acz\'el classification forces $H(t) = \cosh t$, yielding
$J(x) = \tfrac{1}{2}(x + x^{-1}) - 1$.
\end{proof}

The path-weight derivation connects the cost functional to quantum
amplitudes via Feynman's path integral~\cite{Feynman1948}, replacing
the classical action with the RS cost $J$.

\section{Path-Weight Derivation}
\label{sec:path-weight-deriv}

The RS-native derivation proceeds from the path action.

\begin{definition}[Path Action]
\label{def:pathaction}
For a recognition path $\gamma$ on the ledger, the path action
is:
\begin{equation}
  C[\gamma] = \int_\gamma J(r(t))\,dt,
\end{equation}
where $J(x) = \frac{1}{2}(x + x^{-1}) - 1$ is the cost
functional and $r(t)$ is the scale ratio along the path.
\end{definition}

\begin{definition}[Path Weight and Amplitude]
\label{def:weight}
Each path carries a weight $w(\gamma) = e^{-C[\gamma]}$ and an
amplitude:
\begin{equation}
  \psi(\gamma) = e^{-C[\gamma]/2} \cdot e^{i\phi(\gamma)},
\end{equation}
where $\phi(\gamma)$ is the phase accumulated on the 8-tick
cycle.  The modulus--cost bridge is $C = 2A$ where
$A = -\ln|\psi|$, so $|\psi|^2 = e^{-C} = w$.
\end{definition}

\begin{proposition}[Born Rule from Path Weights --- Structural Derivation]
\label{thm:born-path}
The probability of outcome $i$ is:
\begin{equation}
  \boxed{P(i) = |\psi_i|^2
  = \frac{e^{-C_i}}{\sum_j e^{-C_j}}
  = \frac{w_i}{\sum_j w_j},}
\end{equation}
where $\psi_i = \sum_{\gamma \to i} \psi(\gamma)$ is the total
amplitude for paths reaching outcome~$i$.
\end{proposition}

\begin{proof}
The cost of reaching outcome $i$ via a set of paths is
$C_i$, the total path action.  The weight $w_i = e^{-C_i}$
is positive and well-defined.  Normalization requires
$\sum_i P(i) = 1$, so $P(i) = w_i / \sum_j w_j$.  Since
$w_i = |\psi_i|^2$ by the modulus--cost bridge, this is
$P(i) = |\psi_i|^2 / \sum_j |\psi_j|^2$.  For normalized
states ($\sum_j |\psi_j|^2 = 1$), this reduces to
$P(i) = |\psi_i|^2$.
\end{proof}

\begin{remark}[Status of the path-weight derivation]
\label{rem:path-circular}
The path-weight derivation above is structurally motivated but partly
circular: Definition~\ref{def:weight} \emph{defines} $\psi(\gamma) =
e^{-C[\gamma]/2} \cdot e^{i\phi}$, so $|\psi|^2 = e^{-C}$ by
construction.  The modulus-cost bridge $C = 2A$ (with $A = -\ln|\psi|$)
is therefore definitional, not a derived identity.  The non-circular
content of this derivation is that the path-action principle selects
$e^{-C}$ as the natural weight for cost-minimizing dynamics; the
identification of this weight with the quantum probability is the
RS structural postulate.  The uniqueness derivation in
Section~\ref{sec:uniqueness} (Theorem~\ref{thm:born}) provides a
fully independent, non-circular derivation of $P = |c|^2$.
\end{remark}

\begin{remark}
The phase $\phi(\gamma)$ arises from the 8-tick cycle:
position on the $2^D = 8$ clock determines the phase of
each path contribution.  The path weight $w = e^{-C}$
depends only on the cost (magnitude), not on the phase.
This is the origin of $J$-cost phase invariance---phase
carries temporal position information but does not affect
the probability of an outcome.
\end{remark}

\section{The Five Requirements}
\label{sec:requirements}

As an independent cross-check, we derive the Born rule from
five structural requirements on any probability function.

\begin{definition}[Probability Function]
\label{def:prob}
A probability function $P$ assigns to each complex amplitude
$c_i$ a real number $P(c_i)$ representing the probability of
outcome $i$.
\end{definition}

\begin{proposition}[Structural Requirements]
\label{prop:requirements}
The probability function must satisfy:
\begin{enumerate}[label=(\roman*)]
  \item \textbf{Non-negativity}: $P(c_i) \geq 0$ for all
  $c_i$.
  \item \textbf{Normalization}: $\sum_i P(c_i) = 1$ whenever
  $\sum_i |c_i|^2 = 1$.
  \item \textbf{Phase independence}: $P(c_i) = P(|c_i|)$.
  \item \textbf{Multiplicativity}: For independent recognition
  events with amplitudes $c_1$ and $c_2$, the joint probability
  factors: $P(c_1 c_2) = P(c_1)\,P(c_2)$.
  \item \textbf{Continuity}: $P$ is a continuous function of
  the amplitude modulus.
\end{enumerate}
\end{proposition}

\begin{proof}[Justification]
Requirements (i) and (ii) are standard.  Requirement (iii)
follows from the $J$-cost functional: the cost $J(r)$ is
defined on positive reals (scale ratios), so the cost assigned
to an amplitude $c = re^{i\theta}$ is $J(r) = J(|c|)$.  Since
cost determines probability through the path-weight bridge,
and cost is phase-blind, probability is phase-blind.

Requirement (iv) follows from the additivity of path actions:
for independent events on disjoint ledger regions, the total
path action is $C_{12} = C_1 + C_2$.  Therefore
$w_{12} = e^{-C_{12}} = e^{-C_1}\,e^{-C_2} = w_1\,w_2$.
Since $P \propto w = |\psi|^2$, probabilities of independent
events multiply.  At the amplitude level, this means
$P(c_1 c_2) = P(c_1)\,P(c_2)$.

Requirement (v) is physically motivated: measurable quantities
(here, the probability of an outcome as a function of the
amplitude) vary continuously under continuous changes of
experimental parameters.  Mathematically, it is needed to
solve the Cauchy functional equation uniquely (without
continuity, pathological solutions exist that require the
Axiom of Choice).
\end{proof}

\section{Uniqueness of $P = |c|^2$}
\label{sec:uniqueness}

\begin{theorem}[Born Rule from Uniqueness]
\label{thm:born}
The unique function satisfying requirements (i)--(v) is:
\begin{equation}
  \boxed{P(c_i) = |c_i|^2.}
\end{equation}
\end{theorem}

\begin{proof}
By phase independence (iii), $P(c) = f(|c|)$ for some
function $f: [0,\infty) \to [0,\infty)$.

\textbf{Step 1: $f$ is a power law.}
By multiplicativity (iv), for any $r_1, r_2 > 0$:
\begin{equation}
  f(r_1 r_2) = f(r_1)\,f(r_2).
\end{equation}
This is Cauchy's multiplicative functional equation on
$(0,\infty)$~\cite{Aczel1966}.  By requirement~(v)
(continuity), with $f$ non-negative and $f(r) > 0$ for
$r > 0$ (required for non-trivial probabilities), the
unique solution is $f(r) = r^k$ for some constant $k > 0$.

\textbf{Step 2: $k = 2$ from normalization.}
Consider a state $|\psi\rangle = \alpha|a\rangle +
\beta|b\rangle$ with $|\alpha|^2 + |\beta|^2 = 1$.
Normalization (ii) requires
$f(|\alpha|) + f(|\beta|) = 1$.
Setting $|\alpha|^2 = p$ and $|\beta|^2 = 1-p$:
\begin{equation}
  p^{k/2} + (1-p)^{k/2} = 1
  \quad\text{for all } p \in [0,1].
\end{equation}
At $p = 1/2$:
$2 \cdot (1/2)^{k/2} = 1$, so $(1/2)^{k/2} = 1/2$,
giving $k/2 = 1$ and hence $\boldsymbol{k = 2}$.

\emph{Verification}: for $k = 2$, $p^{k/2} + (1-p)^{k/2} =
p + (1-p) = 1$ for all $p \in [0,1]$, confirming that
$k = 2$ satisfies normalization identically. \checkmark

\emph{Note}: The argument uses the normalization convention
$|\alpha|^2 + |\beta|^2 = 1$.  This convention is itself tied to the
Born rule: we are proving that the probability function \emph{consistent}
with this normalization is $P = |c|^2$.  If one normalizes amplitudes
differently (e.g., $|\alpha| + |\beta| = 1$), the analysis changes.
The RS derivation fixes the normalization convention through the path-weight
bridge: $\sum_i e^{-C_i} = 1$ in the vacuum, which is the cost-functional
analog of $\sum_i |c_i|^2 = 1$.

\textbf{Step 3: Exclusion of alternatives.}
For $k = 1$: $\sqrt{p} + \sqrt{1-p} \neq 1$ in general
(at $p = 1/2$: $\sqrt{2} \neq 1$).
For $k = 4$: $p^2 + (1-p)^2 \neq 1$ in general
(at $p = 1/2$: $1/2 \neq 1$).
Only $k = 2$ satisfies normalization for all $p$.
\end{proof}

\begin{remark}
Requirements (i)--(iii) alone do not force $P = |c|^2$.
The functional equation $g(p) + g(1-p) = 1$ (where
$g(p) = f(\sqrt{p})$) has many continuous solutions
beyond $g(p) = p$.  It is requirement (iv)---multiplicativity
from the additivity of path actions---that restricts $f$ to
power laws (given continuity from~(v)), and then normalization
selects $k = 2$.  This is why the $J$-cost structure is
essential: without the cost functional forcing additivity, the
Born rule would not be uniquely determined.
\end{remark}

\section{Interference}
\label{sec:interference}

\begin{lemma}[Interference Identity]
\label{thm:interference}
For amplitudes $c_1, c_2 \in \mathbb{C}$:
\begin{equation}
  |c_1 + c_2|^2 = |c_1|^2 + |c_2|^2
  + 2\,\mathrm{Re}(c_1^* c_2).
\end{equation}
The cross-term $2\,\mathrm{Re}(c_1^* c_2)$ is the interference
term.
\end{lemma}

\begin{proof}
$|c_1 + c_2|^2 = (c_1 + c_2)(c_1 + c_2)^*
= |c_1|^2 + c_1 c_2^* + c_2 c_1^* + |c_2|^2
= |c_1|^2 + |c_2|^2 + 2\,\mathrm{Re}(c_1^* c_2)$.
\end{proof}

\begin{corollary}[Constructive and Destructive Interference]
If $c_1$ and $c_2$ are \emph{in phase} ($\arg c_1 = \arg c_2$),
then $|c_1 + c_2|^2 = (|c_1| + |c_2|)^2$ (constructive
interference).  If they are \emph{out of phase}
($\arg c_1 = \arg c_2 + \pi$), then
$|c_1 + c_2|^2 = (|c_1| - |c_2|)^2$ (destructive interference).
\end{corollary}

\section{Why Not $|\psi|$ or $|\psi|^4$?}
\label{sec:exclusion}

The exclusion of alternative probability rules deserves
emphasis, as it is a common question.

\begin{proposition}[$|\psi|$ Is Excluded]
\label{prop:exclude1}
The rule $P(c) = |c|$ violates normalization.
\end{proposition}

\begin{proof}
For an equal superposition $|\psi\rangle =
\frac{1}{\sqrt{2}}(|a\rangle + |b\rangle)$,
$P(a) + P(b) = \frac{1}{\sqrt{2}} + \frac{1}{\sqrt{2}}
= \sqrt{2} \neq 1$.
\end{proof}

\begin{proposition}[$|\psi|^4$ Is Excluded]
\label{prop:exclude4}
The rule $P(c) = |c|^4$ violates normalization.
\end{proposition}

\begin{proof}
For the same state,
$P(a) + P(b) = (1/\sqrt{2})^4 + (1/\sqrt{2})^4
= 1/4 + 1/4 = 1/2 \neq 1$.
\end{proof}

\begin{corollary}
Any rule $P = |c|^k$ with $k \neq 2$ violates normalization
for at least one state.  The Born rule is the unique survivor.
\end{corollary}

\section{Comparison with Other Derivations}
\label{sec:comparison}

\begin{center}
\renewcommand{\arraystretch}{1.3}
\begin{tabular}{lp{3cm}p{3cm}p{4.5cm}}
\toprule
& \textbf{Gleason} & \textbf{Zurek} & \textbf{RS} \\
\midrule
Starting point & Frame functions on Hilbert space
  & Envariance (entanglement symmetry)
  & $J$-cost path weights \\
Dimension & $\geq 3$ & Any & Any \\
Interpretation & Any & Many-Worlds & Ledger Realism \\
Assumes & Hilbert space structure & Entanglement axioms
  & Cost functional + 8-tick phase \\
Multiplicativity & Implicit & Derived & From path-action additivity \\
Result & $P = |\psi|^2$ & $P = |\psi|^2$ & $P = |\psi|^2$ \\
\bottomrule
\end{tabular}
\end{center}

\begin{remark}
Gleason's theorem requires the Hilbert space dimension $\geq 3$; it
does not apply to qubits (dimension 2).  The RS uniqueness theorem
(Theorem~\ref{thm:born}) applies to \emph{any} dimension, including
qubits, because it relies on cost-functional multiplicativity rather
than properties of projection-valued measures on a Hilbert lattice.
Hardy's~\cite{Hardy2001} derivation from five reasonable axioms and
the RS derivation are structurally complementary: Hardy's most
restrictive axiom (``simplicity'') corresponds to the RS statement
that $J$ is uniquely determined by (A1)--(A3).
Deutsch~\cite{Deutsch1999} and Wallace derive the Born rule
from decision-theoretic axioms within the Everett framework; the
RS derivation replaces the decision-theoretic apparatus with the
$J$-cost multiplicativity, which has a more direct physical
motivation (additivity of path actions).  All approaches arrive
at $P = |\psi|^2$.
\end{remark}

\section{Predictions}
\label{sec:predictions}

\begin{enumerate}
  \item \textbf{Born rule is exact.}  No deviations from
  $P = |\psi|^2$ at any energy scale, including the Planck
  scale.  The rule follows from the structure of the cost
  functional, which does not change with energy.

  \item \textbf{No alternative probability rules.}
  $P = |\psi|$, $|\psi|^4$, and all other power laws are
  excluded (Section~\ref{sec:exclusion}).  Any experiment
  showing $P \neq |\psi|^2$ would falsify the $J$-cost
  framework.

  \item \textbf{Interference is structural.}  All quantum
  interference phenomena follow from the cross-term
  $2\,\mathrm{Re}(c_1^* c_2)$ in the interference identity
  (Section~\ref{sec:interference}), which is a consequence
  of the complex amplitude structure forced by the 8-tick
  cycle.

  \item \textbf{No hidden-variable corrections.}
  The Born rule is not an approximation or an ensemble
  average---it is the exact probability of individual
  outcomes.  RS predicts no sub-quantum deviations of the
  kind sought by nonlinear quantum mechanics proposals.
\end{enumerate}

\section{Conclusion}
\label{sec:conclusion}

The Born rule is derived from RS by two independent routes.
The path-weight derivation (Proposition~\ref{thm:born-path})
proceeds directly from the cost functional: path actions are
additive, weights are multiplicative, and the amplitude--cost
bridge $C = 2A$ gives $P = |\psi|^2$.  The uniqueness derivation (Theorem~\ref{thm:born}) shows that
non-negativity, normalization, phase independence,
multiplicativity, and continuity jointly force $P = |c|^2$ as
the unique probability function.  The
interference identity follows automatically from the complex
amplitude structure.  The Born rule is not postulated but is a
theorem of the cost-geometric framework.

\begin{thebibliography}{99}

\bibitem{RS_Axioms}
J.~Washburn, ``The Algebra of Reality: A Recognition Science Derivation
of Physical Law,'' \emph{Axioms} \textbf{15}(2), 90 (2026).
\texttt{doi:10.3390/axioms15020090}

\bibitem{Born1926}
M.~Born, ``Zur Quantenmechanik der Sto\ss vorg\"ange,''
Z.\ Phys.\ \textbf{37}, 863 (1926).

\bibitem{vN1932}
J.~von~Neumann, \textit{Mathematical Foundations of Quantum
Mechanics} (Princeton University Press, 1932).

\bibitem{Gleason1957}
A.~M.~Gleason, ``Measures on the Closed Subspaces of a Hilbert
Space,'' J.\ Math.\ Mech.\ \textbf{6}, 885 (1957).

\bibitem{Zurek2005}
W.~H.~Zurek, ``Probabilities from Entanglement, Born's Rule
$p_k = |\psi_k|^2$ from Envariance,'' Phys.\ Rev.\ A
\textbf{71}, 052105 (2005).

\bibitem{Feynman1948}
R.~P.~Feynman, ``Space-Time Approach to Non-Relativistic
Quantum Mechanics,'' Rev.\ Mod.\ Phys.\ \textbf{20}, 367
(1948).

\bibitem{Hardy2001}
L.~Hardy, ``Quantum Theory from Five Reasonable Axioms,''
arXiv:quant-ph/0101012 (2001).

\bibitem{Aczel1966}
J.~Acz\'el, \emph{Lectures on Functional Equations and Their
Applications} (Academic Press, 1966).

\bibitem{Deutsch1999}
D.~Deutsch, ``Quantum Theory of Probability and Decisions,''
\emph{Proc.\ R.\ Soc.\ Lond.\ A} \textbf{455}, 3129 (1999).

\end{thebibliography}

\end{document}
